\chapter{Traitement du Risque et Reporting}
\label{chap:traitement}

Une fois les risques identifiés et analysés, l'entreprise doit décider d'une stratégie. C'est ici que la cybersécurité rejoint la gestion de projet et la finance.

\section{Les 4 Options de Traitement du Risque}

La norme ISO 27005 définit quatre options principales pour traiter un risque. Le choix dépend de l'appétence au risque de l'entreprise et du rapport coût/bénéfice.

\subsection{Option 1 : Réduction (Mitigation)}

\begin{boxdef}
\textbf{Réduction :} Mettre en place des mesures de sécurité (techniques ou organisationnelles) pour diminuer la Probabilité et/ou la Gravité du risque. C'est l'option la plus courante.
\end{boxdef}

\begin{boxlizbiz}
\textbf{LiZbiZ (Risque Ransomware) :}
Installer un système de sauvegarde externalisé et cloisonner le réseau (segmentation VLAN).
\textbf{Effet :} La Gravité baisse (on peut restaurer) et la Probabilité baisse (plus difficile à propager).
\end{boxlizbiz}

\subsection{Option 2 : Acceptation}

\begin{boxdef}
\textbf{Acceptation :} Décider consciemment de ne pas traiter le risque. Cela s'applique généralement aux risques faibles ou lorsque le coût du traitement dépasse le coût de l'impact potentiel.
\end{boxdef}

\begin{boxlizbiz}
\textbf{LiZbiZ :} Un risque de panne mineure d'une imprimante du service marketing (impact faible) peut être accepté. On ne mettra pas d'imprimante de secours.
\end{boxlizbiz}

\subsection{Option 3 : Évitement (Suppression)}

\begin{boxdef}
\textbf{Évitement :} Supprimer l'activité ou l'actif qui génère le risque. C'est une décision radicale.
\end{boxdef}

\begin{boxlizbiz}
\textbf{LiZbiZ :} Si la vente d'un produit spécifique (ex: kits de survie) génère trop de fraudes et de risques légaux, LiZbiZ peut décider d'arrêter purement et simplement cette gamme de produits. Le risque disparaît, mais le chiffre d'affaires aussi.
\end{boxlizbiz}

\subsection{Option 4 : Transfert (Partage)}

\begin{boxdef}
\textbf{Transfert :} Déplacer l'impact financier du risque vers un tiers. Cela se fait généralement par le biais d'une assurance ou de l'externalisation.
\end{boxdef}

\begin{boxlizbiz}
\textbf{LiZbiZ :} Souscrire une assurance "Cyber-risque" pour couvrir les frais de rançon et de communication de crise. Ou externaliser l'hébergement du site web à un Cloud Provider (AWS/Azure) pour transférer le risque physique du serveur.
\end{boxlizbiz}

% ------------------------------------------------------------------------------
\section{Le Plan d'Action Sécurité (PAS)}
% ------------------------------------------------------------------------------

\subsection{Définition}

\begin{boxdef}
\textbf{PAS (Plan d'Action Sécurité) :} Document opérationnel qui liste les mesures décidées, les responsables, les budgets et les échéances pour traiter les risques identifiés. C'est la "Feuille de route" du RSSI.
\end{boxdef}

\subsection{Structure type d'un PAS}

Il doit être SMART (Spécifique, Mesurable, Atteignable, Réaliste, Temporel).

\begin{table}[htbp]
\centering
\small
\begin{tabular}{p{4cm}p{2.5cm}p{2.5cm}p{2cm}p{2cm}}
    \toprule
    \textbf{Action} & \textbf{Responsable} & \textbf{Échéance} & \textbf{Budget} & \textbf{Priorité} \\
    \midrule
    Segmentation VLAN (Firewall) & Admin Sys & M+1 & 0€ (Interne) & Haute \\
    Mise à jour Serveur Web & Admin Sys & Immédiat & 0€ & Critique \\
    Formation anti-phishing & RH / RSSI & M+3 & 5000€ & Moyenne \\
    Souscription Assurance & DAF & M+2 & 10k€/an & Moyenne \\
    \bottomrule
\end{tabular}
\caption{Extrait du PAS pour LiZbiZ}
\end{table}

% ------------------------------------------------------------------------------
\section{Reporting et Argumentation Financière}
% ------------------------------------------------------------------------------

Le RSSI doit convaincre le COMEX de financer le PAS. Pour cela, il doit parler le langage de la direction.

\subsection{Le ROI de la Sécurité (ROSPI)}

Il est difficile de calculer un Retour sur Investissement (ROI) classique pour la sécurité, car la sécurité "ne rapporte" pas d'argent direct, elle évite des pertes. On parle plutôt de \textbf{ROSPI} (Return On Security Investment).

\begin{boxdef}
\textbf{ROSPI :} Comparaison entre le coût de la mesure et l'économie de perte potentielle (ALE).
\[ ROSPI = \frac{(Perte\_Évitée - Coût\_Mesure)}{Coût\_Mesure} \]
\end{boxdef}

\subsection{Argumentaire pour le COMEX}

Pour le risque Ransomware (Scénario S1) chez LiZbiZ :

\begin{itemize}
    \item \textbf{Coût de l'incident (Impact) :} 200 000 € (Perte CA + Amende + Image).
    \item \textbf{Probabilité annuelle :} 1 fois par an (Vraisemblance élevée).
    \item \textbf{Coût de la mesure (Segmentation) :} 5 000 € (Jours hommes).
    \item \textbf{Efficacité de la mesure :} Réduit la probabilité de 90\%.
\end{itemize}

\textbf{Calcul simplifié :}
Même si le calcul exact est complexe, l'argument est simple : "Investir 5 000 € maintenant évite un risque de perte de 200 000 € avec une probabilité très forte. C'est une assurance-vie pour l'entreprise."

% ------------------------------------------------------------------------------
\section{Cas LiZbiZ : Présentation au COMEX}
% ------------------------------------------------------------------------------

\subsection{Simulation}

Le DSI et le consultant (étudiant) présentent les résultats au COMEX (le professeur).

\begin{boxlizbiz}
\textbf{Scénario :} Le PDG de LiZbiZ hésite à valider le budget pour la segmentation réseau.

\textbf{Argument du Consultant :}
"Monsieur le PDG, notre analyse EBIOS montre que votre système de facturation est exposé à un risque critique. Actuellement, un pirate peut accéder aux données bancaires des clients depuis le site web public. Le coût de l'aménagement réseau est de 5 jours de travail. Le coût d'une fuite de données est estimé à 200k€ d'amendes RGPD et une perte de crédibilité boursière. Nous recommandons le traitement par \textbf{Réduction} immédiate."
\end{boxlizbiz}

% ------------------------------------------------------------------------------
\section{Exercices de Fin de Module}
% ------------------------------------------------------------------------------

\subsection{Exercice 1 : Décision de Traitement}

Pour chacun des risques suivants chez LiZbiZ, choisissez l'option de traitement la plus adaptée et justifiez :
\begin{enumerate}
    \item Risque de panne électrique générale du petit local serveur (Probabilité Faible, Gravité Haute).
    \item Risque de vol d'un smartphone d'un commercial (Probabilité Moyenne, Gravité Faible si données chiffrées).
    \item Risque juridique lié à l'utilisation d'un logiciel piraté sur un poste isolé.
\end{enumerate}

\subsection{Exercice 2 : Projet Final}

Rédigez une "Fiche de Risque" complète pour le scénario "Vol de données clients via le WiFi Logistique" et proposez un PAS de 3 actions concrètes.
% ==============================================================================
% Fichier : 02_module_risque.tex (Extrait Chapitre 6)
% Description : Evaluation du Module 1
% ==============================================================================